\section{The problem and the result}\label{sec:introduction}

The exact molecular electronic Hamiltonian has unusually rigid input data: its external potential is generated by a list of positive nuclei. A hardness proof for a matrix obtained by projecting onto supplied orbitals does not by itself establish hardness for this operator. Projection provides an upper bound on the ground energy, while a reduction also needs a lower bound that excludes all omitted continuum states. We give such a reduction for the integer-charge, rational-position encoding below.

\subsection{The precise promise problem}

Fix a positive integer $c$. An input consists of distinct positions $R_1,\ldots,R_M\in\Q^3$, positive integers $Z_1,\ldots,Z_M$, an electron number $N\ge1$, and rational thresholds $a<b$. Here $M\ge1$, $N$ is encoded in unary, and the charges are encoded in binary. Every coordinate and threshold is encoded as a signed binary numerator and a positive binary denominator in lowest terms, using a fixed standard self-delimiting encoding. Write $L$ for the total bit length and require $b-a\ge L^{-c}$.

On $\bigwedge^N L^2(\R^3;\C^2)$ consider, in atomic units,
\begin{equation}\label{eq:physical-hamiltonian}
H(R,Z,N)
=-\frac12\sum_{\ell=1}^N\Delta_{x_\ell}
 -\sum_{\ell=1}^N\sum_{\alpha=1}^M
       \frac{Z_\alpha}{|x_\ell-R_\alpha|}
 +\sum_{1\le\ell<k\le N}\frac1{|x_\ell-x_k|}.
\end{equation}
The form domain is
\[
 \mathcal Q_N=H^1((\R^3)^N;(\C^2)^{\otimes N})
       \cap\bigwedge^N L^2(\R^3;\C^2).
\]
Antisymmetry exchanges spatial and spin coordinates together. We use the lower-bounded self-adjoint Coulomb operator associated with the closed form on $\mathcal Q_N$. The three-dimensional Hardy inequality, applied with other coordinates fixed, gives infinitesimal form bounds for the finitely many Coulomb terms and justifies this convention. Set
\[
E_0(R,Z,N)=\inf\Spec H(R,Z,N).
\]
The YES promise is $E_0\le a$, and the NO promise is $E_0\ge b$. No requirement is made outside the promise.

\begin{table}[bp]
\centering
\begin{tabular}{@{}p{.25\linewidth}p{.68\linewidth}@{}}
\toprule
Input or convention & Specification\\
\midrule
Nuclear positions & Pairwise distinct $R_\alpha\in\Q^3$; binary rational coordinates\\
Nuclear charges & Positive integers $Z_\alpha$, encoded in binary\\
Electron number & $N\ge1$, encoded in unary\\
State space & Full spinful antisymmetric continuum space; all spin sectors\\
External data & Nuclear list only; no supplied orbital basis or magnetic field\\
Energy thresholds & Rational $a<b$; YES: $E_0\le a$; NO: $E_0\ge b$\\
Promise precision & $b-a\ge L^{-1}$ in the theorem, for total input length $L$\\
Constructed separation & $b-a\ge1$; a threshold separation, not a spectral gap\\
Energy convention & Clamped-nucleus electronic energy; nuclear repulsion omitted\\
\bottomrule
\end{tabular}
\caption{The input model and the precision attained by the reduction. Rational data use the self-delimiting encoding specified in the text.}
\label{tab:input-model}
\end{table}

There is no supplied basis, electron state, magnetic field, or independently adjustable external potential. The minimization includes every total-spin sector. Nuclear repulsion is omitted because the energy convention is electronic and the nuclei are clamped. The problem places no neutrality, fixed-element, or separate nuclear-separation promise on its inputs. In particular, the threshold gap is a precision requirement; it does not assert a molecular spectral gap, a discrete eigenvalue at the infimum, or an efficiently preparable ground state.

\begin{theorem}\label{thm:main}
The promise problem above is QMA-hard under deterministic classical polynomial-time many-one reductions for $c=1$. The reduction can produce instances with $b-a\ge1$.
\end{theorem}

Thus the exact clamped-nuclei problem is QMA-hard for a fixed inverse-polynomial precision, with the input conventions in \Cref{tab:input-model}. The binary encoding of unrestricted positive nuclear charges is essential to the construction given here. Charges and inverse nuclear separations may have exponentially large values; their descriptions still have polynomial length. The theorem makes no claim of membership in QMA, or of hardness under an additional bound on nuclear charges.

For clarity, QMA-hardness means that every promise problem admitting polynomial-time quantum verification with a polynomial-size quantum witness reduces to this one by a deterministic classical polynomial-time map preserving YES and NO instances. We use a standard QMA-complete Heisenberg problem as the source. The sole complexity-theoretic input is stated precisely in \cref{thm:source}; the subsequent construction and its analytic estimates are proved below.

\subsection{Context and the continuum obstruction}

Electronic-structure complexity has several distinct formulations.
Liu, Christandl, and Verstraete established QMA-completeness of the
finite-mode $N$-representability problem under polynomial-time Turing
reductions, using general two-body fermionic Hamiltonians in the
argument~\cite{LiuChristandlVerstraete2007}. Restricting
the coefficients to those produced by an actual Coulomb potential is a
further requirement. Whitfield, Love, and Aspuru-Guzik explained these
distinctions and the role of magnetic fields in the early hardness
results~\cite[Sections~3.1 and~6.1]{WhitfieldLoveAspuruGuzik2013}.

Schuch and Verstraete proved QMA-completeness for a Hubbard model with
local magnetic fields and developed a continuum realization using designed
scalar and spin-dependent potentials~\cite{SchuchVerstraete2009}. Their
Hubbard-to-Heisenberg route and comparison with the complementary
one-particle space are relevant methodological predecessors. Positive
point nuclei impose a different restriction on the external field; the
continuum comparison needed here is proved in \Cref{sec:continuum}.

O'Gorman, Irani, Whitfield, and Fefferman proved QMA-completeness for electronic structure in a supplied finite basis~\cite[Theorem~1]{OGorman2022}. Their published Appendix~D, Corollary~1, includes positive unit nuclear charges, but the basis remains part of the input; Section~V explicitly raises the complete-space question. Indeed, if $P$ is an orbital projection, the variational inequality
\[
\inf\Spec H\le\min\Spec(PHP|_{\operatorname{Ran}P})
\]
does not transfer a projected NO-instance lower bound to $H$. A continuum hardness proof must exclude lower states outside that projection.

The Heisenberg hardness theorem of Cubitt, Montanaro, and Piddock~\cite{CMP2018,CMP2019} supplies a field-free square-lattice source with signed polynomially bounded interactions. Its efficient spatially sparse construction is relevant: generic simulation of a dense interaction graph would not by itself give the required polynomial interaction scales. We use the singlet mediator of Piddock and Montanaro~\cite{PiddockMontanaro2017} to obtain a positive Heisenberg model, and then use repulsive Hubbard superexchange~\cite{Anderson1959}. The finite calculations below supply explicit scalar shifts, errors, rationalization, and the geometric control required by the nuclear construction.

Schleich et al.~\cite{Schleich2026} study chemically motivated preparation
and dynamics through assumptions on accessible chemical processes. Those
assumptions define a different computational task from deciding every
promised spectral-infimum instance considered here.

The companion article \emph{QMA-hardness of continuum Coulomb energy with
unit nuclear charges}~\cite[Theorem~1]{UnitCoulomb2026} proves hardness
with every $Z_\alpha=1$. Its construction uses positive continuous charge
density, polynomial confinement, and equal-mass cubature. That theorem
implies hardness of the larger binary-charge input class, but it does not
use the binary-charge theorem proved here. Conversely, splitting our
exponentially large charges into individual unit nuclei would not give
polynomial output size. The present proof retains a distinct quantitative
contribution: certified exponential-accuracy eigenvalue computation and
relative tunneling calibration for widely separated hydrogenic wells.

Binding questions are also distinct from the decision problem. For the
same kinetic normalization, the companion ionization theorem
\cite[Theorem~1.1]{UniformExcessCharge2026} states that strict binding
$E_N<E_{N-1}$ implies $N\le Z_{\rm tot}+CM$ for a universal $C$, where
$Z_{\rm tot}=\sum_\alpha Z_\alpha$ and $E_0=0$ in that particle-number
notation. It assumes distinct nuclei and charges at least one, with no
separation condition. Our reduction imposes no strict-binding promise,
so that bound supplies neither the continuum exclusion estimate nor a
ground-state existence assertion. It is a scope comparison only.

\subsection{Mechanism of the reduction}

We first replace each signed Heisenberg bond by positive interactions with a strongly bound ancillary singlet. A repulsive Hubbard model at half filling then realizes these positive exchanges through virtual double occupancy: its low-energy singly occupied states carry the spins, and the effective exchange is proportional to the square of the hopping strength. The continuum construction must therefore realize prescribed small positive hoppings together with a fixed onsite repulsion. The continuum comparison retains all configurations and spin sectors in the selected orbitals; single occupancy is obtained in the finite Hubbard analysis, not imposed on admissible continuum states.

The construction uses a large integer scale $Z_*\asymp e^D$, where $D$ is a fixed polynomial in the source size. In dilated coordinates $y=Z_*x$, each site contains a nucleus of charge close to $Z_*$ and a weaker positive secondary nucleus. Their normalized one-electron states are close to hydrogen ground states. A secondary displacement tunes each isolated-well energy of the dilated operator to $-1/2$ with exponentially small error. The effective model uses the shifted physical energy divided by $Z_*$. In these units, hopping amplitudes are $Z_*$ times their dilated one-electron values, while the onsite electron repulsion approaches the fixed constant $5/8$.

Three quantitative ingredients make the continuum and bit-complexity requirements compatible.

\begin{enumerate}[label=(\roman*)]
\item The sites are placed so that small displacements can independently change every required bond length to first order, with a quantitative bound on the displacement size. A relative tail estimate, proved with capped remote fields and comparison barriers, converts these changes into controlled hopping ratios. Absolute proximity to a hydrogen orbital would not give the required relative precision.
\item A graded polynomial Ritz space evaluates a one-well eigenvalue to exponential accuracy in polynomial time. Uniform analytic estimates, an explicit rational mass-matrix lower bound, and certified eigenvalue enclosures make the procedure valid even for skinny cells and ill-conditioned intermediate bases. The unknown eigenfunction is used only to prove approximation, not as an input to the algorithm.
\item Localization yields a fixed gap on the entire one-particle complement. A kinetic-trace and Hardy estimate controls Coulomb coupling on every fermionic superposition in the selected modes. Completing a square against the complementary gap gives a lower bound for the full many-electron spectral infimum.
\end{enumerate}

These arguments provide quantitative tools for constructions involving exponentially weak tunneling and a continuum representation. Their uniformity is part of the proof: errors must remain small after multiplication by $Z_*$, and all required precision must cost only polynomially many bits. The numerical construction combines local polynomial approximation through a partition of unity, in the sense of Melenk and Babu\v{s}ka~\cite{MelenkBabuska1996}, with separate rational conditioning and integration bounds. The complement estimate uses spatial localization~\cite{Simon1983} and a quadratic elimination argument whose finite-dimensional analogue is second-order perturbation theory~\cite{BravyiDiVincenzoLoss2011}. These methods are credited at their points of use; their needed quantitative forms are proved here.

The order of construction and the order of justification deserve separate attention. We first prescribe the finite Hubbard hoppings and compute the primary centers from positive integrals of explicit capped Coulomb fields. These surrogate fields do not involve any unknown orbital or secondary displacement. Only after rounding the primary centers and fixing the integer charges do we tune the secondary nuclei. The surrogate replacement bounds are uniform throughout the tuning cube, and the relative orbital estimates apply to every tuned output. Finally, the continuum lower bound transfers the Hubbard NO promise to arbitrary wavefunctions, and an affine change of energy transfers both promises to rational physical thresholds.

\Cref{sec:finite} develops the finite models and the contact layout. \Cref{sec:wells} proves the isolated-well estimates and the eigenvalue oracle. \Cref{sec:hopping} computes the positions and calibrates hopping. \Cref{sec:continuum} proves the full continuum comparison. \Cref{sec:completion} chooses all parameters, gives the classical reduction in its actual order of execution, and verifies the final thresholds and encoding.

Throughout, constants denoted by $C,c>0$ are independent of the large scale parameters and may change between estimates. A polynomial bound always means a fixed polynomial with uniform coefficients and degree. Exponentially small bounds are stated explicitly when they are later amplified. All operator inequalities involving unbounded Hamiltonians are quadratic-form inequalities unless their bounded compressions are specified.
